% Conditional enclosure incorporated in v0.8; scope remains conditional.
% Numerical error budget does not remove the inherited analytical hypotheses.
\subsection{A conditional enclosure for the monopole response}
\label{nl:conditional-monopole}
For the real dipole tangent at T2, we assembled a complete error budget
for the $L=0$ second-harmonic forced problem. Here $L$ labels the
radiated angular sector, whereas the underlying localized linear mode
has $l=1$. The calculation uses the exact-mode Jost normalization,
with the evanescent reduced component satisfying
$\lim_{r\to\infty}e^{\kappa_c r}v(r)=1$.
It therefore has a different amplitude normalization from the numerical
power table in Section~\ref{nl:dipole-response}.

The error identity includes the inner and outer source differences,
operator differences, the residual of the fixed approximate response,
its actual matching jumps and origin contribution, and the exterior
boundary and coefficient-extraction errors. The exterior potential and
source are included through infinity. No measured small linewidth or
comparison between two numerical resolutions substitutes for these terms.

The following inherited hypotheses specify the exact problem to which
the numerical error budget applies; they are not new conclusions of
the scalar summation:
\begin{enumerate}
\item The T2 existence enclosure of Section~\ref{sec:certified-existence}
supplies one common exact parameter tuple
$z_*=(a_*,c_*,\rho_*,\omega_*)$ in its certified dyadic box, a positive
regular spherical profile $f_*$, and a nonzero real regular dipole mode.
All profile, mode and parameter-error bounds refer to this same tuple.
Neither uniqueness of the tuple nor the claim that every point of the
box is an exact root is assumed.

\item At that exact tuple, the decaying fundamental Jost solution belongs
to the span of the two regular origin solutions. This is the inherited
regular/Jost identification from the T2 Wronskian conditions, not an
inference from overlap of numerical residual intervals. Its reduced
components are $O(r^2)$ at the origin and obey the same normalization
$e^{\kappa_c r}v(r)\to1$, with
$\kappa_c=\sqrt{1-(\omega_*-\rho_*)^2}$.
The background and mode satisfy the analytic infinite-tail bounds:
for $r\ge36$ one may use
$|f_*(r)|\le(97/8)e^{-0.495r}/r$ and
$|u(r)|+|v(r)|\le1.097e^{-0.287r}$.

\item The exact quadratic $L=0$ source is the complete product source
of the real standing dipole in Eqs.~\eqref{nl:quadratic}--\eqref{nl:forced}.
The fixed nominal mode, source, response and comparison coefficient
use the same amplitude convention: the nominal unit-$L^2$ mode is divided
by the exact rational
$s_0=2558188928468177/2251799813685248$, and quadratic quantities by $s_0^2$.
This does not identify $s_0$ with the norm of the exact mode.
The computational carrier conjugation is retained throughout.

\item For the second harmonic, the exact reduced operator is
$-\partial_r^2-Q_*^2+V_*(r)$ with real symmetric, locally bounded
$V_*$ and
$Q_*=\operatorname{diag}(q_+,q_-)$,
$q_\pm=\sqrt{(\omega_*\pm2\rho_*)^2-1}$.
Both channels satisfy $5/2<q_\pm<23/5$.
The exterior Volterra bounds hold at these exact frequencies,
with the analytic potential and source tails retained through infinity.
At $R=64$, the exterior potential satisfies
$\int_R^\infty\|V_*(r)\|_2\,dr/(5/2)<1/2$ and the source is integrable.
Thus the exact outgoing reduction includes both its homogeneous maps
and its inhomogeneous source corrections. Together with regular
Dirichlet data at the origin, the positive outgoing matrix current
gives the unique forced solution and adjoint extraction problem.
This argument concerns two open second-harmonic channels, not the
fundamental embedded mode.

\item The analytic and validated input bounds used in the error identity
enclose the entire core and exterior, including exact-to-nominal
frequency and map changes. In Euclidean vector and induced matrix norms,
the inherited bounds
$\int_0^\infty\|V_*(r)\|_2\,dr<25/4$ and the Volterra estimates give
the adjoint amplitude weights $\|Z\|_2<26/5$ and $\|Z'\|_2<14$.
Source and operator differences, response residuals, origin and matching
jumps, exterior source terms, and coefficient-extraction rounding are
all bounded in this same convention, without dropping a tail or counting
one twice. The directed-arithmetic implementation and its interval
libraries remain part of the computational trust assumptions.
\end{enumerate}
Under these inherited hypotheses, the computed common coefficient-error
bound is

\begin{equation}
 \|C_* - C_{\mathrm{num}}\|_2
 \leq \mathcal E < 0.001023269269845280.
\end{equation}
The frozen comparison coefficient in the second computational channel
has the lower bound
\begin{equation}
 |C_{\mathrm{num},2}| > 0.001582591541472706.
\end{equation}
The directed arithmetic establishes, before rounding for display,
\begin{equation}
 |C_{*,2}| > 0.000559322271627427.
\end{equation}
This channel has laboratory frequency $\omega_*-2\rho_*<0$.
The solver uses a negative carrier and conjugates its second computational
component to recover that physical sideband. Conjugating the entire field
to the positive-carrier convention used here gives
$C_{-,\mathrm{paper}}=C_{*,2}$; all these operations preserve the modulus.
The first channel remains undecided by this common error bound; that
failure does not imply its coefficient vanishes.

This result supplies the nonzero-coefficient premise of the conditional
smooth-continuation argument for a real standing dipole. Its extension
to a complex dipole polarization is proportional to
$\boldsymbol p^T\boldsymbol p$ and therefore does not cover circular
polarization. It gives neither a nonlinear lifetime nor a dynamical
instability theorem. In particular, the bound is not a certified value
for any power coefficient in the preceding approximate-input table.
